% !PRINTED 673
\weberpage{673}

If $\mathfrak{P}$ is any point and $\varrho$ a function which has order 1 at this
point, then following § 185 we set

\weq{2}{\begin{aligned}
\eta_1 &= e_1\varrho^m + \sigma_1\varrho^{m+1},\\
\eta_2 &= e_2\varrho^m + \sigma_2\varrho^{m+1},\\
&\text{\quad$\cdot\ \cdot\ \cdot\ \cdot\ \cdot\ \cdot\ \cdot\ \cdot\ \cdot\ \cdot$}\\
\eta_s &= e_s\varrho^m + \sigma_s\varrho^{m+1},
\end{aligned}}

in which $e_1, e_2, \ldots, e_s$ are constants that do not all vanish,
and $\sigma_1, \sigma_2, \ldots, \sigma_s$ are functions that are finite at $\mathfrak{P}$;
$m$ is the lowest order number at $\mathfrak{P}$ that occurs among the functions $\eta_i$.

The totality of the functions

\weq{3}{\eta = c_1\eta_1 + c_2\eta_2 + \cdots + c_s\eta_s}

with constants $c_1, c_2, \ldots, c_s$ is called a \gesperrt{family}. We denote it by

\weq{4}{(\eta_1, \eta_2, \ldots, \eta_s),}

and call $\eta_1, \eta_2, \ldots, \eta_s$ a \gesperrt{basis of the family}. The order
of $\eta$ at $\mathfrak{P}$ is not smaller than $m$ (it is greater when
$e_1c_1 + e_2c_2 + \cdots + e_sc_s = 0$). But $\eta$ can become infinite at no
point $\mathfrak{P}$ which is not contained in $\mathfrak{A}$, and
also not to a higher order than that to which the point $\mathfrak{P}$ enters into $\mathfrak{A}$.

We can therefore set:

\weq{5}{\eta = \frac{\mathfrak{A}'}{\mathfrak{A}},}

in which $\mathfrak{A}'$ likewise belongs to the class $A$.

If instead of $\mathfrak{A}$ one takes another polygon $\mathfrak{B}$ of the class $A$
and sets

\weq{6}{\xi = \frac{\mathfrak{A}}{\mathfrak{B}},}

then from 1. it follows that:

\weq{7}{\begin{aligned}
\xi\eta_1 &= \eta_1' = \frac{\mathfrak{A}_1}{\mathfrak{B}},\\
\xi\eta_2 &= \eta_2' = \frac{\mathfrak{A}_2}{\mathfrak{B}},\\
&\text{\quad$\cdot\ \cdot\ \cdot\ \cdot\ \cdot\ \cdot\ \cdot\ \cdot$}\\
\xi\eta_s &= \eta_s' = \frac{\mathfrak{A}_s}{\mathfrak{B}}
\end{aligned}}

and

\weq{8}{\xi\eta = \eta' = c_1\eta_1' + c_2\eta_2' + \cdots + c_s\eta_s' = \frac{\mathfrak{A}'}{\mathfrak{B}}.}

Every polygon $\mathfrak{A}'$ generated by the denominator $\mathfrak{A}$ and a system
of constants $c_1, c_2, \ldots, c_s$ is therefore also generated by every
%CONT>>
